% Section 14.8, from lectures/L14/appendix/b_exercises.md.

\section{Exercises}
\label{c14:sec:exercises}

\begin{aside}
These exercises consolidate the chapter. Design your own \file{tx_shift_reg.vhd} from the
specification in \secref{c14:sec:iface} to \ref{c14:sec:trace}; the testbench handed out, run per
appendix~\ref{app:sim}, is the check that it behaves correctly.
\end{aside}

\exercise{\code{tx_shift_reg}}{VHDL}
\label{c14:ex:txsr}
Write your own \file{tx_shift_reg.vhd} from the chapter. Verify it:

\begin{shell}
cd controller
ghdl -a --std=93 can_def.vhd tx_shift_reg.vhd tx_shift_reg_tb.vhd
ghdl -e --std=93 tx_shift_reg_tb
ghdl -r --std=93 tx_shift_reg_tb --assert-level=error
\end{shell}

If a check fails, read the assertion message; it names exactly which shift it concerns and what was
expected.

\exercise{Trace a stuffed sequence by hand, then in simulation}{Simulation}
\label{c14:ex:trace}
\xpart{a} Trace by hand \code{tx_shift_reg} loaded with \code{data = "11100000"},
\code{bit_count = 8}, starting at a group boundary where the previous group's last bit was \code{0}
with a run length of 3 (that is, \code{consecutive = 3} and \code{last_bit = '0'} going into the
group). Write out, bit by bit: what is presented at \code{load}, and then at every subsequent
\code{shift}, with every inserted stuff bit marked and stating which shift raises \code{done}.

\xpart{b} Confirm your hand trace in simulation: write a short testbench (or reuse the structure of
\file{tx_shift_reg_tb.vhd}) that preloads the same starting conditions with run length 3, then loads
\code{"11100000"} and at every shift \code{assert}s \code{tx_bit}, \code{stuff} and \code{done}
against the values you predicted (or \code{report}s them so that you can read them). Does its
pass-or-fail outcome agree with your hand trace?

\exercise{The one-bit group that still owes a stuff bit}{Simulation}
\label{c14:ex:onebit}
\Secref{c14:sec:timing}'s last point names an edge case without working it through: \emph{if
\code{bit_count = 1} and that single bit triggers a stuff bit, \code{done} still waits for the stuff
bit to be sent first.} That is not hypothetical; \code{can_controller} loads one-bit groups (SOF) and
four-bit groups (the ID's lower half plus RTR), so a group really can end mid-run.

Set it up exactly as tests 2 and 3 in \code{tx_shift_reg_tb} do, with no reset in between: load
\code{"00001111"} with \code{bit_count = 8} and shift it out completely, so that the run counter
stands at four \code{'1'}s in a row. Then load a \textbf{one-bit group} whose single bit is also
\code{'1'} (\code{data = "1-------"}, \code{bit_count = 1}).

\xpart{a} Trace what happens, presentation by presentation: what \code{load} puts on \code{tx_bit},
what the run counter reaches, and then what every following \code{shift} does. Mark which
presentation carries the stuff bit and which raises \code{done}.

\xpart{b} How many \code{shift} pulses does this one-bit group consume in total, and how many would
it have consumed if the run counter had instead stood at zero? Say where the difference comes from.

\xpart{c} Suppose \code{done} had been allowed to pulse as soon as the group's single real bit was
presented, with no regard for the pending stuff bit. \code{can_controller} would then load the next
group during the next bit period. Describe exactly what the receiving node would then see on the
wire, and which of \code{rx_shift_reg}'s outputs would eventually report it (\chapref{c15:ch}).

\xpart{d} Confirm your walkthrough in simulation, by extending the pattern from test 2 and test 3 in
\file{tx_shift_reg_tb.vhd} with this one-bit group as a fourth case.

\exercise{Where \code{track_run} is called from}{Simulation}
\label{c14:ex:trackrun}
\Secref{c14:sub:tracker} puts the run counter in a procedure and is explicit about where that
procedure may be called from: from a branch placing a bit, never from the process's top level. Both
halves of that rule are worth breaking deliberately, in a scratch copy of your module.

\xpart{a} Move the call \code{track_run(data(7));} out of the \code{load} branch, so that \code{load}
latches the group and presents its first bit but the counter never sees that bit. Predict what the
run counter then shows for the rest of the frame, and which of the testbench's three tests fails
first. Run it.

\xpart{b} Now put a single \code{track_run(tx_bit);} at the top of the
\code{elsif (rising_edge(clock))} branch, above the defaults, and remove all three calls inside the
branches. \code{tx_bit} is an \code{out} port, so start by saying why that does not even compile, and
then repeat the experiment with an internal copy of the transmitted bit in its place. With
\code{TICKS_PER_BIT} set to 50, how many times does the counter now run per CAN bit, and what does
\code{consecutive} do on the second clock edge of the very first bit period?

\xpart{c} \Secref{c14:sub:tracker} also says that the procedure has to be declared \emph{inside} the
process rather than in the architecture. Move it out to the architecture's declarative part,
unchanged, and analyse the file. It does not compile; read the errors, one per signal assignment in
the procedure, before you go on. Then make an architecture-level version that \emph{does} compile,
with identical behaviour and a passing testbench. Compare the two call sites and say what the second
version costs, and which new question it raises that the inside-the-process version cannot raise.
